% answers.tex —— 2009 年·普通高等学校招生考试·山东卷·文科 · 参考答案与详细解析片段
% 由 tools/generate_projects.py 依来源总账自 2009/2009_数学真题_山东文_试题.md + 2009/2009_数学真题_山东文_解析正文.md + 2009/2009_数学真题_山东文_答案要点.md 生成，勿手改（改源 MD 后重跑）


\begin{gktip}
固定候选 PDF 仅含题面，无答案、解析和评分标准，原卷亦未标分值；以下为非官方编者独立推导，官方答案字段为空。题面文字层只作定位，分数、根号、上下标与五幅图均以 600 dpi 原页复核。
第 17 题原页给出「$f(x)=2\sin x\cos^2\frac x2+\cos x\sin\varphi-\sin x$ 在 $x=\pi$ 处取最小值」，化简后 $f'(\pi)=1\neq0$，条件与函数不相容，故该题记 U 并登记问题，详见该题解析。
\end{gktip}

\begin{answer}
\gksolhead{1}{}

\textbf{题目}：

集合 $A=\{0,2,a\}$，$B=\{1,a^2\}$．若 $A\cup B=\{0,1,2,4,16\}$，则 $a$ 的值为（　）

（A）$0$

（B）$1$

（C）$2$

（D）$4$

\Needspace{4\baselineskip}
\textbf{解析}

由 $A\cup B=\{0,1,2,a,a^2\}=\{0,1,2,4,16\}$ 得 $\{a,a^2\}=\{4,16\}$。

若 $a=4$，则 $a^2=16$，并集恰为 $\{0,1,2,4,16\}$，符合；若 $a=16$，则 $a^2=256$，并集会出现 $256$，不合。

故 $a=4$，选 D。

\gkanswer{答案：D。}

\end{answer}
\begin{answer}
\gksolhead{2}{}

\textbf{题目}：

复数 $\dfrac{3-\mathrm i}{1-\mathrm i}$ 等于（　）

（A）$1+2\mathrm i$

（B）$1-2\mathrm i$

（C）$2+\mathrm i$

（D）$2-\mathrm i$

\Needspace{4\baselineskip}
\textbf{解析}

\nopagebreak[4]
\[\frac{3-\mathrm i}{1-\mathrm i}
=\frac{(3-\mathrm i)(1+\mathrm i)}{(1-\mathrm i)(1+\mathrm i)}
=\frac{3+3\mathrm i-\mathrm i-\mathrm i^2}{2}
=\frac{4+2\mathrm i}{2}=2+\mathrm i,\]

选 C。

\gkanswer{答案：C。}

\end{answer}
\begin{answer}
\gksolhead{3}{}

\textbf{题目}：

将函数 $y=\sin2x$ 的图象向左平移 $\dfrac{\pi}{4}$ 个单位，再向上平移 $1$ 个单位，所得图象的函数解析式是（　）

（A）$y=\cos2x$

（B）$y=2\cos^2x$

（C）$y=1+\sin\left(2x+\dfrac{\pi}{4}\right)$

（D）$y=2\sin^2x$

\Needspace{4\baselineskip}
\textbf{解析}

\Needspace{5\baselineskip}
向左平移 $\frac\pi4$ 得 $y=\sin2\left(x+\frac\pi4\right)=\sin\left(2x+\frac\pi2\right)=\cos2x$，再向上平移 $1$ 得

\nopagebreak[4]
\[y=1+\cos2x=2\cos^2x,\]

选 B。

\gkanswer{答案：B。}

\end{answer}
\begin{answer}
\gksolhead{4}{}

\textbf{题目}：

一空间几何体的三视图如图所示，则该几何体的体积为（　）

（A）$2\pi+2\sqrt3$

（B）$4\pi+2\sqrt3$

（C）$2\pi+\dfrac{2\sqrt3}{3}$

（D）$4\pi+\dfrac{2\sqrt3}{3}$
\nopagebreak[4]
\begin{center}% 2009 年山东卷文科第 4 题：三视图（圆柱与其上四棱锥）。
% 来源 PDF：普通高考/2009/2009山东文.pdf，SHA256 1807901a21da57e3d45382294189b0c8e795722b6879842319ae75428639fad6
% 原页：第 1 页（A3 三栏）第 1 栏中部，印刷页 547。
% 重绘依据：正视图为 2×2 矩形上方接顶角三角形（斜棱标 2），侧视图同形；
%   俯视图为半径 1 的圆内接正方形（对角线为两条互相垂直的直径）；尺寸按原图标注。
\begin{tikzpicture}[scale=0.8, semithick, every node/.style={inner sep=1.3pt, font=\small}]
  % 正视图
  \begin{scope}[shift={(0,0)}]
    \draw (0,0) rectangle (2,2);
    \draw (0,2) -- (1,3.0) -- (2,2);
    \node[anchor=east] at (-0.12,1) {$2$};
    \node[anchor=north] at (1,-0.12) {$2$};
    \node[anchor=south east] at (0.42,2.55) {$2$};
    \node[anchor=north] at (1,-0.42) {正（主）视图};
  \end{scope}
  % 侧视图
  \begin{scope}[shift={(4.2,0)}]
    \draw (0,0) rectangle (2,2);
    \draw (0,2) -- (1,3.0) -- (2,2);
    \node[anchor=north] at (1,-0.12) {$2$};
    \node[anchor=south east] at (0.42,2.55) {$2$};
    \node[anchor=north] at (1,-0.42) {侧（左）视图};
  \end{scope}
  % 俯视图
  \begin{scope}[shift={(1.0,-4.6)}]
    \draw (1,1) circle (1);
    \draw (0,1) -- (2,1) -- (1,2) -- (1,0) -- cycle;
    \node[anchor=north] at (1,-0.18) {俯视图};
  \end{scope}
\end{tikzpicture}
\end{center}
\nopagebreak[4]
\Needspace{4\baselineskip}
\textbf{解析}

由三视图：下部是底面半径 $1$、高 $2$ 的圆柱；上部是四棱锥，其底面为圆柱底面的内接正方形（俯视图中圆内接的正方形，对角线长为 $2$），侧棱长为 $2$。

\Needspace{5\baselineskip}
圆柱体积

\nopagebreak[4]
\[V_1=\pi\cdot1^2\cdot2=2\pi .\]

\Needspace{5\baselineskip}
内接正方形对角线为 $2$，边长 $\sqrt2$，中心到顶点距离 $\frac{\sqrt2\cdot\sqrt2}{2}=1$，故四棱锥的高

\nopagebreak[4]
\[h=\sqrt{2^2-1^2}=\sqrt3,\]

\nopagebreak[4]
\[V_2=\frac13\cdot(\sqrt2)^2\cdot\sqrt3=\frac{2\sqrt3}{3}.\]

\nopagebreak[4]
\[V=V_1+V_2=2\pi+\frac{2\sqrt3}{3},\]

选 C。

\gkanswer{答案：C。}

\end{answer}
\begin{answer}
\gksolhead{5}{}

\textbf{题目}：

在 $\mathbf R$ 上定义运算 $\odot$：$a\odot b=ab+2a+b$，则满足 $x\odot(x-2)<0$ 的实数 $x$ 的取值范围为（　）

（A）$(0,2)$

（B）$(-2,1)$

（C）$(-\infty,-2)\cup(1,+\infty)$

（D）$(-1,2)$

\Needspace{4\baselineskip}
\textbf{解析}

\nopagebreak[4]
\[x\odot(x-2)=x(x-2)+2x+(x-2)=x^2-2x+2x+x-2=x^2+x-2 .\]

\nopagebreak[4]
\[x^2+x-2<0\iff (x+2)(x-1)<0\iff -2<x<1,\]

选 B。

\gkanswer{答案：B。}

\end{answer}
\begin{answer}
\gksolhead{6}{}

\textbf{题目}：

函数 $y=\dfrac{\mathrm e^x+\mathrm e^{-x}}{\mathrm e^x-\mathrm e^{-x}}$ 的图象大致为（　）
\nopagebreak[4]
\begin{center}% 2009 年山东卷文科第 6 题：y=(e^x+e^{-x})/(e^x-e^{-x}) 的四个候选图象。
% 来源 PDF：普通高考/2009/2009山东文.pdf，SHA256 1807901a21da57e3d45382294189b0c8e795722b6879842319ae75428639fad6
% 原页：第 1 页（A3 三栏）第 2 栏上部，印刷页 547。
% 重绘依据：(A) 双曲余切：右支由 +∞ 递减趋于 1，左支由 −1 递减到 −∞，虚线为 y=±1；
%   (B) 在 y 轴处有跳跃（(0,1) 与 (0,−1) 为空心点），整体递增；
%   (C) 过原点、以 y 轴为竖直切线、两侧递增（形如 y=x^{1/3}）；
%   (D) 关于 y 轴对称、在原点取 0、两侧递增（形如 y=|x|^{1/3}）。
\begin{tikzpicture}[semithick, every node/.style={inner sep=1.1pt, font=\small}]
  % (A)
  \begin{scope}[shift={(0,0)}, x=0.62cm, y=0.62cm]
    \draw[dashed] (-2.0,1) -- (2.4,1);
    \draw[dashed] (-2.0,-1) -- (2.4,-1);
    \draw[-{Stealth}] (-2.1,0) -- (2.5,0) node[right] {$x$};
    \draw[-{Stealth}] (0,-2.2) -- (0,2.2) node[above] {$y$};
    \begin{scope}
      \clip (-2.1,-2.2) rectangle (2.5,2.2);
      \draw[domain=0.14:2.3,samples=140,smooth,variable=\x]
        plot (\x, {(exp(\x)+exp(-\x))/(exp(\x)-exp(-\x))});
      \draw[domain=-2.3:-0.14,samples=140,smooth,variable=\x]
        plot (\x, {(exp(\x)+exp(-\x))/(exp(\x)-exp(-\x))});
    \end{scope}
    \node[anchor=north] at (0.02,-0.10) {$O$};
    \node[anchor=south east] at (-0.06,1) {$1$};
    \node[anchor=north east] at (-0.06,-1) {$-1$};
    \node at (0.6,-2.85) {(A)};
  \end{scope}
  % (B)
  \begin{scope}[shift={(5.0,0)}, x=0.62cm, y=0.62cm]
    \draw[-{Stealth}] (-2.1,0) -- (2.5,0) node[right] {$x$};
    \draw[-{Stealth}] (0,-2.2) -- (0,2.2) node[above] {$y$};
    \begin{scope}
      \clip (-2.1,-2.2) rectangle (2.5,2.2);
      \draw[domain=-2.2:-0.001,samples=100,smooth,variable=\x]
        plot (\x, {1+0.55*\x+0.10*\x*\x});
      \draw[domain=0.001:2.2,samples=100,smooth,variable=\x]
        plot (\x, {-1+0.55*\x+0.10*\x*\x});
    \end{scope}
    \draw[fill=white] (0,1) circle (1.4pt);
    \draw[fill=white] (0,-1) circle (1.4pt);
    \node[anchor=north] at (0.02,-0.10) {$O$};
    \node[anchor=west] at (0.08,1) {$1$};
    \node at (0.6,-2.85) {(B)};
  \end{scope}
  % (C)
  \begin{scope}[shift={(0,-5.4)}, x=0.62cm, y=0.62cm]
    \draw[-{Stealth}] (-2.1,0) -- (2.5,0) node[right] {$x$};
    \draw[-{Stealth}] (0,-2.2) -- (0,2.2) node[above] {$y$};
    \begin{scope}
      \clip (-2.1,-2.2) rectangle (2.5,2.2);
      \draw[domain=-2.2:2.2,samples=160,smooth,variable=\x]
        plot (\x, {sign(\x)*pow(abs(\x),1/3)*1.25});
    \end{scope}
    \node[anchor=north] at (0.02,-0.10) {$O$};
    \node[anchor=south west] at (1,-0.02) {$1$};
    \node at (0.6,-2.85) {(C)};
  \end{scope}
  % (D)
  \begin{scope}[shift={(5.0,-5.4)}, x=0.62cm, y=0.62cm]
    \draw[-{Stealth}] (-2.1,0) -- (2.5,0) node[right] {$x$};
    \draw[-{Stealth}] (0,-2.2) -- (0,2.2) node[above] {$y$};
    \begin{scope}
      \clip (-2.1,-2.2) rectangle (2.5,2.2);
      \draw[domain=0.001:2.2,samples=100,smooth,variable=\x]
        plot (\x, {pow(\x,1/3)*1.25});
      \draw[domain=-2.2:-0.001,samples=100,smooth,variable=\x]
        plot (\x, {pow(-\x,1/3)*1.25});
    \end{scope}
    \node[anchor=north] at (0.02,-0.10) {$O$};
    \node[anchor=south east] at (-0.05,1.25) {$1$};
    \node at (0.6,-2.85) {(D)};
  \end{scope}
\end{tikzpicture}
\end{center}
\nopagebreak[4]
\Needspace{4\baselineskip}
\textbf{解析}

\nopagebreak[4]
\[f(x)=\frac{\mathrm e^x+\mathrm e^{-x}}{\mathrm e^x-\mathrm e^{-x}},\qquad x\neq0,\]

\Needspace{5\baselineskip}
且

\nopagebreak[4]
\[f(-x)=\frac{\mathrm e^{-x}+\mathrm e^{x}}{\mathrm e^{-x}-\mathrm e^{x}}=-f(x),\]

为奇函数。

\Needspace{5\baselineskip}
$x>0$ 时分子分母同除以 $\mathrm e^x$：

\nopagebreak[4]
\[f(x)=\frac{1+\mathrm e^{-2x}}{1-\mathrm e^{-2x}},\]

随 $x$ 增大而减小，$x\to0^+$ 时 $f\to+\infty$，$x\to+\infty$ 时 $f\to1$。由奇对称，$x<0$ 时 $f(x)<-1$、递减，$x\to-\infty$ 时 $f\to-1$。

故图象以 $y$ 轴为竖直渐近线、以 $y=\pm1$ 为水平渐近线，选 A。

\gkanswer{答案：A。}

\end{answer}
\begin{answer}
\gksolhead{7}{}

\textbf{题目}：

定义在 $\mathbf R$ 上的函数 $f(x)$ 满足 $f(x)=\begin{cases}\log_2(4-x),&x\leqslant0,\\ f(x-1)-f(x-2),&x>0,\end{cases}$ 则 $f(3)$ 的值为（　）

（A）$-1$

（B）$-2$

（C）$1$

（D）$2$

\Needspace{4\baselineskip}
\textbf{解析}

\nopagebreak[4]
\[f(0)=\log_2(4-0)=2,\qquad f(-1)=\log_2(4+1)=\log_2 5 .\]

\Needspace{5\baselineskip}
由 $f(x)=f(x-1)-f(x-2)$ 依次取 $x=1,2,3$：

\nopagebreak[4]
\[f(1)=f(0)-f(-1)=2-\log_2 5,\]

\nopagebreak[4]
\[f(2)=f(1)-f(0)=-\log_2 5,\]

\nopagebreak[4]
\[f(3)=f(2)-f(1)=-\log_2 5-(2-\log_2 5)=-2,\]

选 B。

\gkanswer{答案：B。}

\end{answer}
\begin{answer}
\gksolhead{8}{}

\textbf{题目}：

设 $P$ 是 $\triangle ABC$ 所在平面内的一点，$\overrightarrow{BC}+\overrightarrow{BA}=2\overrightarrow{BP}$，则（　）

（A）$\overrightarrow{PA}+\overrightarrow{PB}=\mathbf 0$

（B）$\overrightarrow{PB}+\overrightarrow{PC}=\mathbf 0$

（C）$\overrightarrow{PC}+\overrightarrow{PA}=\mathbf 0$

（D）$\overrightarrow{PA}+\overrightarrow{PB}+\overrightarrow{PC}=\mathbf 0$

\Needspace{4\baselineskip}
\textbf{解析}

\Needspace{5\baselineskip}
由

\nopagebreak[4]
\[\overrightarrow{BC}+\overrightarrow{BA}=2\overrightarrow{BP}\]

\Needspace{5\baselineskip}
得

\nopagebreak[4]
\[\overrightarrow{BP}=\frac{\overrightarrow{BA}+\overrightarrow{BC}}{2},\]

\Needspace{5\baselineskip}
即 $P$ 是线段 $AC$ 的中点。于是 $\overrightarrow{PA}$ 与 $\overrightarrow{PC}$ 大小相等、方向相反，

\nopagebreak[4]
\[\overrightarrow{PC}+\overrightarrow{PA}=\mathbf 0,\]

选 C。

\gkanswer{答案：C。}

\end{answer}
\begin{answer}
\gksolhead{9}{}

\textbf{题目}：

已知 $\alpha,\beta$ 表示两个不同的平面，$m$ 为平面 $\alpha$ 内的一条直线，则「$\alpha\perp\beta$」是「$m\perp\beta$」的（　）

（A）充分不必要条件

（B）必要不充分条件

（C）充要条件

（D）既不充分也不必要条件

\Needspace{4\baselineskip}
\textbf{解析}

若 $m\perp\beta$，又 $m\subset\alpha$，则由面面垂直的判定定理得 $\alpha\perp\beta$，故「$\alpha\perp\beta$」是必要条件。

反之，$\alpha\perp\beta$ 时，平面 $\alpha$ 内只有垂直于两平面交线的直线才垂直于 $\beta$；取 $m$ 为 $\alpha$ 内与交线平行的直线，则 $m\parallel\beta$，不垂直于 $\beta$。故不充分。

选 B。

\gkanswer{答案：B。}

\end{answer}
\begin{answer}
\gksolhead{10}{}

\textbf{题目}：

设斜率为 $2$ 的直线 $l$ 过抛物线 $y^2=ax$（$a\neq0$）的焦点 $F$，且和 $y$ 轴交于点 $A$．若 $\triangle OAF$（$O$ 为坐标原点）的面积为 $4$，则抛物线方程为（　）

（A）$y^2=\pm4x$

（B）$y^2=\pm8x$

（C）$y^2=4x$

（D）$y^2=8x$

\Needspace{4\baselineskip}
\textbf{解析}

$y^2=ax$ 的焦点为 $F\left(\frac a4,0\right)$（$a\neq0$）。

\Needspace{5\baselineskip}
斜率为 $2$ 且过 $F$ 的直线：

\nopagebreak[4]
\[y=2\left(x-\frac a4\right).\]

\Needspace{5\baselineskip}
令 $x=0$ 得 $A\left(0,-\frac a2\right)$。于是

\nopagebreak[4]
\[S_{\triangle OAF}=\frac12|OF|\cdot|OA|=\frac12\cdot\left|\frac a4\right|\cdot\left|\frac a2\right|=\frac{a^2}{16}.\]

\Needspace{5\baselineskip}
由 $S=4$ 得 $a^2=64$，即 $a=\pm8$。故抛物线为

\nopagebreak[4]
\[y^2=\pm8x,\]

选 B。

\gkanswer{答案：B。}

\end{answer}
\begin{answer}
\gksolhead{11}{}

\textbf{题目}：

在区间 $\left[-\dfrac{\pi}{2},\dfrac{\pi}{2}\right]$ 上随机取一个数 $x$，$\cos x$ 的值介于 $0$ 到 $\dfrac12$ 之间的概率为（　）

（A）$\dfrac13$

（B）$\dfrac2\pi$

（C）$\dfrac12$

（D）$\dfrac23$

\Needspace{4\baselineskip}
\textbf{解析}

\Needspace{5\baselineskip}
在 $\left[-\frac\pi2,\frac\pi2\right]$ 上 $\cos x\geqslant0$，且

\nopagebreak[4]
\[0<\cos x<\frac12\iff \frac\pi3<|x|\leqslant\frac\pi2 .\]

\Needspace{5\baselineskip}
故 $x$ 落在 $\left[-\frac\pi2,-\frac\pi3\right)$ 与 $\left(\frac\pi3,\frac\pi2\right]$ 内，两段总长

\nopagebreak[4]
\[\frac\pi2-\frac\pi3+\frac\pi2-\frac\pi3=\frac\pi3 .\]

\Needspace{5\baselineskip}
区间总长为 $\pi$，所求概率

\nopagebreak[4]
\[P=\frac{\pi/3}{\pi}=\frac13,\]

选 A。

\gkanswer{答案：A。}

\end{answer}
\begin{answer}
\gksolhead{12}{}

\textbf{题目}：

已知定义在 $\mathbf R$ 上的奇函数 $f(x)$ 满足 $f(x-4)=-f(x)$，且在区间 $[0,2]$ 上是增函数，则（　）

（A）$f(-25)<f(11)<f(80)$

（B）$f(80)<f(11)<f(-25)$

（C）$f(11)<f(80)<f(-25)$

（D）$f(-25)<f(80)<f(11)$

\Needspace{4\baselineskip}
\textbf{解析}

\Needspace{5\baselineskip}
由 $f(x-4)=-f(x)$，把 $x$ 换成 $x+4$ 得 $f(x)=-f(x+4)$，即 $f(x+4)=-f(x)$。于是

\nopagebreak[4]
\[f(x+8)=-f(x+4)=f(x),\]

\Needspace{5\baselineskip}
故 $8$ 是 $f$ 的一个周期。又 $f$ 为奇函数，$f(0)=0$。

\nopagebreak[4]
\[f(80)=f(0)=0,\]

\nopagebreak[4]
\[f(11)=f(11-8)=f(3),\]

\Needspace{5\baselineskip}
而由 $f(3)=-f(-1)=f(1)$（用到 $f$ 为奇函数）得 $f(11)=f(1)$；

\nopagebreak[4]
\[f(-25)=f(-25+32)=f(7)=-f(3)=-f(1).\]

\Needspace{5\baselineskip}
$f$ 在 $[0,2]$ 上递增且 $f(0)=0$，故 $f(1)>0$。于是

\nopagebreak[4]
\[f(-25)=-f(1)<0=f(80)<f(1)=f(11),\]

即 $f(-25)<f(80)<f(11)$，选 D。

\gkanswer{答案：D。}

\end{answer}
\begin{answer}
\gksolhead{13}{}

\textbf{题目}：

在等差数列 $\{a_n\}$ 中，$a_3=7$，$a_5=a_2+6$，则 $a_6=$\gkblank{}．

\Needspace{4\baselineskip}
\textbf{解析}

\Needspace{5\baselineskip}
设公差为 $d$，则

\nopagebreak[4]
\[a_5-a_2=3d=6\ \Longrightarrow\ d=2 .\]

\Needspace{5\baselineskip}
故

\nopagebreak[4]
\[a_6=a_3+3d=7+6=13 .\]

\gkanswer{答案：$13$。}

\end{answer}
\begin{answer}
\gksolhead{14}{}

\textbf{题目}：

若函数 $f(x)=a^x-x-a$（$a>0$ 且 $a\neq1$）有两个零点，则实数 $a$ 的取值范围是\gkblank{}．

\Needspace{4\baselineskip}
\textbf{解析}

\nopagebreak[4]
\[f'(x)=a^x\ln a-1 .\]

\begin{itemize}
\item 若 $0<a<1$，则 $\ln a<0$，$a^x\ln a<0$，故 $f'(x)<0$ 恒成立，$f$ 在 $\mathbf R$ 上严格递减，至多一个零点。
\item 若 $a>1$，则 $f''(x)=a^x(\ln a)^2>0$，$f$ 为凸函数，且
\end{itemize}

\nopagebreak[4]
\[f'(x)=0\iff x=x_0=\log_a\frac1{\ln a}.\]

\Needspace{5\baselineskip}
此时 $x\to-\infty$ 时 $f(x)=a^x-x-a\to+\infty$，$x\to+\infty$ 时 $f\to+\infty$，故 $f$ 有两个零点当且仅当 $f(x_0)<0$。

\nopagebreak[4]
\[f(0)=1-0-a=1-a<0,\]

故 $f(x_0)\leqslant f(0)<0$，条件自动满足。

综上，$a$ 的取值范围是 $(1,+\infty)$。

\gkanswer{答案：$(1,+\infty)$。}

\end{answer}
\begin{answer}
\gksolhead{15}{}

\textbf{题目}：

执行如图的程序框图，输出的 $T=$\gkblank{}．
\nopagebreak[4]
\begin{center}% 2009 年山东卷文科第 15 题：程序框图，输出 T 的值。
% 来源 PDF：普通高考/2009/2009山东文.pdf，SHA256 1807901a21da57e3d45382294189b0c8e795722b6879842319ae75428639fad6
% 原页：第 1 页（A3 三栏）第 3 栏上部，印刷页 547。
% 重绘依据：开始 → S=0,T=0,n=0 → 判断 T>S：是（右）输出 T 并结束；否（下）S=S+5 → n=n+2 → T=T+n，
%   再回到判断。
\begin{tikzpicture}[scale=0.9, semithick, >=Stealth,
  node distance=5mm, every node/.style={inner sep=3pt, font=\small}]
  \node (start) [draw, rounded corners=4pt, minimum width=1.9cm, minimum height=0.62cm] {开始};
  \node (init) [draw, minimum width=3.6cm, minimum height=0.62cm, below=of start] {$S=0,\ T=0,\ n=0$};
  \node (cond) [draw, diamond, aspect=2.2, minimum width=3.0cm, minimum height=0.95cm, below=of init] {$T>S$};
  \node (s) [draw, minimum width=2.2cm, minimum height=0.62cm, below=of cond] {$S=S+5$};
  \node (n) [draw, minimum width=2.2cm, minimum height=0.62cm, below=of s] {$n=n+2$};
  \node (t) [draw, minimum width=2.2cm, minimum height=0.62cm, below=of n] {$T=T+n$};
  \node (out) [draw, trapezium, trapezium left angle=70, trapezium right angle=110,
    minimum width=2.2cm, minimum height=0.62cm, right=1.6cm of cond] {输出 $T$};
  \node (end) [draw, rounded corners=4pt, minimum width=1.9cm, minimum height=0.62cm, below=of out] {结束};
  \draw[->] (start) -- (init);
  \draw[->] (init) -- (cond);
  \draw[->] (cond) -- node[right, pos=0.55] {否} (s);
  \draw[->] (s) -- (n);
  \draw[->] (n) -- (t);
  \draw[->] (t.west) -- ++(-1.3,0) |- (cond.west);
  \draw[->] (cond) -- node[above, pos=0.45] {是} (out);
  \draw[->] (out) -- (end);
\end{tikzpicture}
\end{center}
\nopagebreak[4]
\Needspace{4\baselineskip}
\textbf{解析}

初始 $S=T=n=0$。每次判断 $T>S$：若成立则输出 $T$ 并结束，否则执行 $S\leftarrow S+5$、$n\leftarrow n+2$、$T\leftarrow T+n$。

\begin{center}
\begin{tabular}{lllll}
\toprule
轮次 & $T>S$ & $S$ & $n$ & $T$ \\
\midrule
1 & $0>0$ 否 & $5$ & $2$ & $2$ \\
2 & $2>5$ 否 & $10$ & $4$ & $6$ \\
3 & $6>10$ 否 & $15$ & $6$ & $12$ \\
4 & $12>15$ 否 & $20$ & $8$ & $20$ \\
5 & $20>20$ 否 & $25$ & $10$ & $30$ \\
6 & $30>25$ 是 & — & — & 输出 $30$ \\
\bottomrule
\end{tabular}
\end{center}

故输出 $T=30$。

\gkanswer{答案：$30$。}

\end{answer}
\begin{answer}
\gksolhead{16}{}

\textbf{题目}：

某公司租赁甲、乙两种设备生产 $A,B$ 两类产品，甲种设备每天能生产 $A$ 类产品 $5$ 件和 $B$ 类产品 $10$ 件，乙种设备每天能生产 $A$ 类产品 $6$ 件和 $B$ 类产品 $20$ 件．已知设备甲每天的租赁费为 $200$ 元，设备乙每天的租赁费为 $300$ 元．现该公司至少要生产 $A$ 类产品 $50$ 件，$B$ 类产品 $140$ 件，所需租赁费最少为\gkblank{}元．

\Needspace{4\baselineskip}
\textbf{解析}

\Needspace{5\baselineskip}
设租用甲种设备 $x$ 天、乙种设备 $y$ 天，则

\nopagebreak[4]
\[5x+6y\geqslant50,\qquad 10x+20y\geqslant140,\]

\Needspace{5\baselineskip}
即

\nopagebreak[4]
\[5x+6y\geqslant50,\qquad x+2y\geqslant14,\qquad x,y\geqslant0 .\]

\Needspace{5\baselineskip}
目标函数为

\nopagebreak[4]
\[z=200x+300y .\]

可行域的三个顶点：

\begin{itemize}
\item 两直线交点：由 $x=14-2y$ 代入 $5x+6y=50$ 得 $70-10y+6y=50$，$y=5$、$x=4$，即 $(4,5)$；
\item $x=0$ 与 $x+2y=14$：$y$ 需同时满足 $6y\geqslant50$，得 $\left(0,\frac{25}{3}\right)$；
\item $y=0$ 与 $x+2y=14$：$(14,0)$。
\end{itemize}

\Needspace{5\baselineskip}
目标值

\nopagebreak[4]
\[z(4,5)=2300,\qquad z\left(0,\frac{25}{3}\right)=2500,\qquad z(14,0)=2800 .\]

故最少租赁费为 $2300$ 元。

\gkanswer{答案：$2300$。}

\end{answer}
\begin{answer}
\gksolhead{17}{}

\textbf{题目}：

已知函数 $f(x)=2\sin x\cos^2\dfrac x2+\cos x\sin\varphi-\sin x$（$0<\varphi<\pi$）在 $x=\pi$ 处取最小值．

（1）求 $\varphi$ 的值；

（2）在 $\triangle ABC$ 中，$a,b,c$ 分别是角 $A,B,C$ 的对边．已知 $a=1$，$b=\sqrt2$，$f(A)=\dfrac{\sqrt3}{2}$，求角 $C$．

\Needspace{4\baselineskip}
\textbf{解析}

\Needspace{5\baselineskip}
先化简：

\nopagebreak[4]
\[2\sin x\cos^2\frac x2=\sin x(1+\cos x),\]

\nopagebreak[4]
\[f(x)=\sin x+\sin x\cos x+\cos x\sin\varphi-\sin x=\cos x(\sin x+\sin\varphi).\]

\Needspace{5\baselineskip}
求导：

\nopagebreak[4]
\[f'(x)=\cos2x-\sin x\sin\varphi .\]

\Needspace{5\baselineskip}
在 $x=\pi$ 处

\nopagebreak[4]
\[f'(\pi)=\cos2\pi-\sin\pi\sin\varphi=1-0=1\neq0 .\]

\Needspace{5\baselineskip}
若 $x=\pi$ 是极小值点，必须 $f'(\pi)=0$，矛盾。进一步，取 $0<\varepsilon\ll1$：

\nopagebreak[4]
\[f(\pi-\varepsilon)-f(\pi)=\cos(\pi-\varepsilon)(\sin(\pi-\varepsilon)+\sin\varphi)+\sin\varphi .\]

\Needspace{5\baselineskip}
由 $\cos(\pi-\varepsilon)=-\cos\varepsilon\approx-1+\frac{\varepsilon^2}{2}$、$\sin(\pi-\varepsilon)=\sin\varepsilon\approx\varepsilon$ 得

\nopagebreak[4]
\[f(\pi-\varepsilon)\approx\left(-1+\frac{\varepsilon^2}{2}\right)(\varepsilon+\sin\varphi)+\sin\varphi\approx-\sin\varphi-\varepsilon+\sin\varphi=-\varepsilon<0,\]

即 $f(\pi-\varepsilon)<f(\pi)$，故 $\pi$ 不可能是最小值点。

因此按原页字面第 (1) 问无解，第 (2) 问的 $\varphi$ 也无从确定。

（若有人只由 $f(\pi)=-\sin\varphi$「取最小值」而认定 $\sin\varphi=1$，得 $\varphi=\frac\pi2$；但这个取值既与 $f'(\pi)=0$ 的必要条件矛盾，又会使第 (2) 问无解：此时 $f(A)=\cos A(1+\sin A)=\frac{\sqrt3}{2}$ 要求 $A\approx1.1$ rad，而 $a=1$、$b=\sqrt2$ 时由正弦定理 $\sin B=\sqrt2\sin A>1$，无三角形。故本库不采用该值。）

本库不代拟答案，第 17 题整体记 U，并登记 ISSUE-1482（转录缺陷）与 ISSUE-1483（条件不相容）。

\gkanswer{答案：U（原页条件与所给函数不相容，无法确定 $\varphi$）。}

\end{answer}
\begin{answer}
\gksolhead{18}{}

\textbf{题目}：

如图，在直四棱柱 $ABCD-A_1B_1C_1D_1$ 中，底面 $ABCD$ 为等腰梯形，$AB\parallel CD$，$AB=4$，$BC=CD=2$，$AA_1=2$，$E$、$E_1$ 分别是棱 $AD$、$AA_1$ 的中点．
\nopagebreak[4]
\begin{center}% 2009 年山东卷文科第 18 题：直四棱柱 ABCD-A1B1C1D1，底面为等腰梯形。
% 来源 PDF：普通高考/2009/2009山东文.pdf，SHA256 1807901a21da57e3d45382294189b0c8e795722b6879842319ae75428639fad6
% 原页：第 2 页（A3 三栏）第 1 栏上部，印刷页 548。
% 重绘依据：底面 ABCD 按斜投影画成等腰梯形（AB 为下底、CD 为上底），A 左下、B 右下、C 右上、D 左上；
%   A1B1C1D1 为对应上底；E 在 AD 上、E1 在 AA1 上、F 在 AB 上；
%   被遮挡的棱 A1D1、D1C1、DD1 与辅助线画虚线。
\begin{tikzpicture}[scale=0.95, semithick, every node/.style={inner sep=1.3pt}]
  \coordinate (A) at (0,0);
  \coordinate (B) at (4.0,0);
  \coordinate (C) at (2.9,1.15);
  \coordinate (D) at (1.1,1.15);
  \coordinate (A1) at (0,2.2);
  \coordinate (B1) at (4.0,2.2);
  \coordinate (C1) at (2.9,3.35);
  \coordinate (D1) at (1.1,3.35);
  \coordinate (E) at ($(A)!0.5!(D)$);
  \coordinate (E1) at ($(A)!0.5!(A1)$);
  \coordinate (F) at ($(A)!0.5!(B)$);
  % 底面
  \draw (A) -- (B) -- (C) -- (D) -- cycle;
  % 上底（部分被遮挡）
  \draw (A1) -- (B1) -- (C1);
  \draw[dashed] (C1) -- (D1) -- (A1);
  % 侧棱
  \draw (A) -- (A1);
  \draw (B) -- (B1);
  \draw (C) -- (C1);
  \draw[dashed] (D) -- (D1);
  % 点
  \fill (E) circle (1.1pt);
  \fill (E1) circle (1.1pt);
  \fill (F) circle (1.1pt);
  \node[anchor=north west] at (A) {$A$};
  \node[anchor=north east] at (B) {$B$};
  \node[anchor=west] at (C) {$C$};
  \node[anchor=south east] at (D) {$D$};
  \node[anchor=east] at (A1) {$A_1$};
  \node[anchor=south east] at (B1) {$B_1$};
  \node[anchor=south west] at (C1) {$C_1$};
  \node[anchor=east] at (D1) {$D_1$};
  \node[anchor=north west] at (E) {$E$};
  \node[anchor=east] at (E1) {$E_1$};
  \node[anchor=north] at ([yshift=-2pt]F) {$F$};
\end{tikzpicture}
\end{center}
\nopagebreak[4]
（1）设 $F$ 是棱 $AB$ 的中点，证明：直线 $EE_1\parallel$ 平面 $FCC_1$；

（2）证明：平面 $D_1AC\perp$ 平面 $BB_1C_1C$．

\Needspace{4\baselineskip}
\textbf{解析}

\Needspace{5\baselineskip}
取坐标系：$A(0,0,0)$、$B(4,0,0)$。底面等腰梯形中 $CD=2$、$BC=2$、$AB=4$，故

\nopagebreak[4]
\[D(1,\sqrt3,0),\qquad C(3,\sqrt3,0).\]

\Needspace{5\baselineskip}
直四棱柱高 $AA_1=2$：

\nopagebreak[4]
\[A_1(0,0,2),\quad B_1(4,0,2),\quad C_1(3,\sqrt3,2),\quad D_1(1,\sqrt3,2).\]

\Needspace{5\baselineskip}
$E$、$E_1$、$F$ 分别为 $AD$、$AA_1$、$AB$ 的中点：

\nopagebreak[4]
\[E\left(\frac12,\frac{\sqrt3}{2},0\right),\quad E_1(0,0,1),\quad F(2,0,0).\]

\Needspace{5\baselineskip}
（1）

\nopagebreak[4]
\[\overrightarrow{EE_1}=\left(-\frac12,-\frac{\sqrt3}{2},1\right).\]

\Needspace{5\baselineskip}
平面 $FCC_1$ 内

\nopagebreak[4]
\[\overrightarrow{FC}=(1,\sqrt3,0),\qquad \overrightarrow{CC_1}=(0,0,2),\]

\Needspace{5\baselineskip}
法向量

\nopagebreak[4]
\[\boldsymbol n=\overrightarrow{FC}\times\overrightarrow{CC_1}
=(2\sqrt3,-2,0)\parallel(\sqrt3,-1,0).\]

\nopagebreak[4]
\[\overrightarrow{EE_1}\cdot\boldsymbol n=-\frac{\sqrt3}{2}+\frac{\sqrt3}{2}+0=0 .\]

\Needspace{5\baselineskip}
又平面 $FCC_1$ 的方程为 $\sqrt3(x-2)-y=0$，代入 $E$：

\nopagebreak[4]
\[\sqrt3\left(\frac12-2\right)-\frac{\sqrt3}{2}=-\frac{3\sqrt3}{2}-\frac{\sqrt3}{2}=-2\sqrt3\neq0,\]

\Needspace{5\baselineskip}
即 $E\notin$ 平面 $FCC_1$。故

\nopagebreak[4]
\[EE_1\parallel\text{\text{平面} }FCC_1 .\]

\Needspace{5\baselineskip}
（2）

\nopagebreak[4]
\[\overrightarrow{AD_1}=(1,\sqrt3,2),\qquad \overrightarrow{AC}=(3,\sqrt3,0),\]

\nopagebreak[4]
\[\boldsymbol n_1=\overrightarrow{AD_1}\times\overrightarrow{AC}
=\left(\sqrt3\cdot0-2\cdot\sqrt3,\ 2\cdot3-1\cdot0,\ 1\cdot\sqrt3-\sqrt3\cdot3\right)
=(-2\sqrt3,6,-2\sqrt3),\]

\Needspace{5\baselineskip}
取 $(-\sqrt3,3,-\sqrt3)$。

\nopagebreak[4]
\[\overrightarrow{BB_1}=(0,0,2),\qquad \overrightarrow{BC}=(-1,\sqrt3,0),\]

\nopagebreak[4]
\[\boldsymbol n_2=\overrightarrow{BB_1}\times\overrightarrow{BC}
=(0-2\sqrt3,\ -2-0,\ 0)=(-2\sqrt3,-2,0),\]

\Needspace{5\baselineskip}
取 $(-\sqrt3,-1,0)$。

\nopagebreak[4]
\[\boldsymbol n_1\cdot\boldsymbol n_2=3-3+0=0,\]

故平面 $D_1AC\perp$ 平面 $BB_1C_1C$。

\gkanswer{答案：（1）（2）证明见解析。}

\end{answer}
\begin{answer}
\gksolhead{19}{}

\textbf{题目}：

一汽车厂生产 $A,B,C$ 三类轿车，每类轿车均有舒适型和标准型两种型号，某月的产量如下表（单位：辆）：
\nopagebreak[4]
\begin{center}% 2009 年山东卷文科第 19 题：三类轿车两型号的产量表。
% 来源 PDF：普通高考/2009/2009山东文.pdf，SHA256 1807901a21da57e3d45382294189b0c8e795722b6879842319ae75428639fad6
% 原页：第 2 页（A3 三栏）第 1 栏下部，印刷页 548。
% 重绘依据：表头为「车型」与轿车 A、B、C，两行为舒适型、标准型；A 列 100/300，B 列 150/450，
%   C 列舒适型为 z、标准型为 600；单位「辆」随题干给出。
\begin{tikzpicture}[scale=1.0, semithick, every node/.style={inner sep=3pt, font=\small}]
  \draw (0,0) rectangle (7.2,-1.6);
  \draw (0,-0.8) -- (7.2,-0.8);
  \draw (1.8,0) -- (1.8,-1.6);
  \draw (3.6,0) -- (3.6,-1.6);
  \draw (5.4,0) -- (5.4,-1.6);
  \node at (0.9,-0.4) {车型};
  \node at (2.7,-0.4) {轿车 $A$};
  \node at (4.5,-0.4) {轿车 $B$};
  \node at (6.3,-0.4) {轿车 $C$};
  \node at (0.9,-1.2) {舒适型};
  \node at (2.7,-1.2) {$100$};
  \node at (4.5,-1.2) {$150$};
  \node at (6.3,-1.2) {$z$};
  \draw (0,-2.4) -- (7.2,-2.4);
  \draw (0,0) -- (0,-2.4);
  \draw (1.8,0) -- (1.8,-2.4);
  \draw (3.6,0) -- (3.6,-2.4);
  \draw (5.4,0) -- (5.4,-2.4);
  \draw (7.2,0) -- (7.2,-2.4);
  \node at (0.9,-2.0) {标准型};
  \node at (2.7,-2.0) {$300$};
  \node at (4.5,-2.0) {$450$};
  \node at (6.3,-2.0) {$600$};
\end{tikzpicture}
\end{center}
\nopagebreak[4]
按类型分层抽样的方法在这个月生产的轿车中抽取 $50$ 辆，其中有 $A$ 类轿车 $10$ 辆．

（1）求 $z$ 的值；

（2）用分层抽样的方法在 $C$ 类轿车中抽取一个容量为 $5$ 的样本．将该样本看成一个总体，从中任取 $2$ 辆，求至少有 $1$ 辆舒适型轿车的概率；

（3）用随机抽样的方法从 $B$ 类舒适型轿车中抽取 $8$ 辆，经检测它们的得分如下：$9.4$，$8.6$，$9.2$，$9.6$，$8.7$，$9.3$，$9.0$，$8.2$．把这 $8$ 辆轿车的得分看成一个总体，从中任取一个数，求该数与样本平均数之差的绝对值不超过 $0.5$ 的概率．

\Needspace{4\baselineskip}
\textbf{解析}

\Needspace{5\baselineskip}
（1）$A$ 类轿车共 $100+300=400$ 辆，抽取 $10$ 辆，抽样比

\nopagebreak[4]
\[\frac{10}{400}=\frac1{40}.\]

\Needspace{5\baselineskip}
总产量为 $50\div\frac1{40}=2000$ 辆，故

\nopagebreak[4]
\[(100+300)+(150+450)+(z+600)=2000\ \Longrightarrow\ z=400 .\]

\Needspace{5\baselineskip}
（2）$C$ 类轿车中舒适型 $400$ 辆、标准型 $600$ 辆，比例 $2:3$。容量 $5$ 的样本中舒适型 $2$ 辆、标准型 $3$ 辆。

\nopagebreak[4]
\[\mathrm C_5^2=10,\qquad \text{\text{两辆都是标准型的情况}}\ \mathrm C_3^2=3,\]

\nopagebreak[4]
\[P(\text{\text{至少有} }1\text{ \text{辆舒适型}})=1-\frac3{10}=\frac7{10}.\]

\Needspace{5\baselineskip}
（3）样本平均数

\nopagebreak[4]
\[\bar x=\frac{9.4+8.6+9.2+9.6+8.7+9.3+9.0+8.2}{8}=\frac{72.0}{8}=9.0 .\]

\Needspace{5\baselineskip}
$|x_i-9.0|\leqslant0.5\iff 8.5\leqslant x_i\leqslant9.5$，满足的数据为 $9.4,8.6,9.2,8.7,9.3,9.0$ 共 $6$ 个，故

\nopagebreak[4]
\[P=\frac68=\frac34 .\]

\gkanswer{答案：（1）$z=400$；（2）$\frac7{10}$；（3）$\frac34$。}

\end{answer}
\begin{answer}
\gksolhead{20}{}

\textbf{题目}：

等比数列 $\{a_n\}$ 的前 $n$ 项和为 $S_n$．已知对任意的 $n\in\mathbf N^*$，点 $(n,S_n)$ 均在函数 $y=b^x+r$（$b>0$ 且 $b\neq1$，$b,r$ 为常数）的图象上．

（1）求 $r$ 的值；

（2）当 $b=2$ 时，记 $b_n=\dfrac{n+1}{4a_n}$（$n\in\mathbf N^*$），求数列 $\{b_n\}$ 的前 $n$ 项和 $T_n$．

\Needspace{4\baselineskip}
\textbf{解析}

（1）由题意 $S_n=b^n+r$。

\Needspace{5\baselineskip}
$n=1$ 时 $a_1=S_1=b+r$；$n\geqslant2$ 时

\nopagebreak[4]
\[a_n=S_n-S_{n-1}=b^n-b^{n-1}=(b-1)b^{n-1}.\]

\Needspace{5\baselineskip}
因 $\{a_n\}$ 是等比数列，$a_1$ 必须等于 $\frac{a_2}{b}=b-1$（$b\neq0$、$b\neq1$），故

\nopagebreak[4]
\[b+r=b-1\ \Longrightarrow\ r=-1 .\]

验证：此时 $a_1=b-1$、$a_n=(b-1)b^{n-1}$，公比为 $b$，确为等比数列。

\Needspace{5\baselineskip}
（2）$b=2$ 时 $a_n=2^{n-1}$，

\nopagebreak[4]
\[b_n=\frac{n+1}{4\cdot2^{n-1}}=\frac{n+1}{2^{n+1}} .\]

\nopagebreak[4]
\[T_n=\sum_{k=1}^{n}\frac{k+1}{2^{k+1}}=\frac14\sum_{k=1}^{n}(k+1)\left(\frac12\right)^{k-1}.\]

\Needspace{5\baselineskip}
由

\nopagebreak[4]
\[\sum_{k=1}^{n}k x^{k-1}=\frac{1-(n+1)x^n+nx^{n+1}}{(1-x)^2},\qquad
\sum_{k=1}^{n}x^{k-1}=\frac{1-x^n}{1-x},\]

\Needspace{5\baselineskip}
取 $x=\frac12$：

\nopagebreak[4]
\[\sum_{k=1}^{n}(k+1)\left(\frac12\right)^{k-1}
=\frac{1-\frac{n+1}{2^n}+\frac{n}{2^{n+1}}}{1/4}+\frac{1-\frac1{2^n}}{1/2}
=4\left(1-\frac{n+1}{2^n}+\frac{n}{2^{n+1}}\right)+2\left(1-\frac1{2^n}\right).\]

\Needspace{5\baselineskip}
化简：

\nopagebreak[4]
\[=6-\frac{4(n+1)}{2^n}+\frac{4n}{2^{n+1}}-\frac{2}{2^n}
=6-\frac{4n+4+2-2n}{2^n}=6-\frac{2n+6}{2^n}.\]

\Needspace{5\baselineskip}
故

\nopagebreak[4]
\[T_n=\frac14\left(6-\frac{2n+6}{2^n}\right)=\frac32-\frac{n+3}{2^{n+1}} .\]

\gkanswer{答案：（1）$r=-1$；（2）$T_n=\frac32-\frac{n+3}{2^{n+1}}$。}

\end{answer}
\begin{answer}
\gksolhead{21}{}

\textbf{题目}：

已知函数 $f(x)=\dfrac13ax^3+bx^2+x+3$，其中 $a\neq0$．

（1）当 $a,b$ 满足什么条件时，$f(x)$ 取得极值？

（2）已知 $a>0$，且 $f(x)$ 在区间 $(0,1]$ 上单调递增，试用 $a$ 表示 $b$ 的取值范围．

\Needspace{4\baselineskip}
\textbf{解析}

\Needspace{5\baselineskip}
（1）

\nopagebreak[4]
\[f'(x)=ax^2+2bx+1\qquad(a\neq0).\]

\Needspace{5\baselineskip}
$f$ 取得极值当且仅当 $f'=0$ 有两个不相等的实根：

\nopagebreak[4]
\[\Delta=4b^2-4a>0\iff b^2>a .\]

\Needspace{5\baselineskip}
（2）$a>0$ 时 $f'(x)=ax^2+2bx+1\geqslant0$ 对一切 $x\in(0,1]$ 成立，即

\nopagebreak[4]
\[2b\geqslant-\left(ax+\frac1x\right)\qquad(0<x\leqslant1).\]

记 $g(x)=ax+\frac1x$，则 $g'(x)=a-\frac1{x^2}$，$g$ 在 $x=\frac1{\sqrt a}$ 处取最小值（若该点在 $(0,1]$ 内）。

\begin{itemize}
\item $0<a<1$ 时 $\frac1{\sqrt a}>1$，$g$ 在 $(0,1]$ 上递减，最小值为 $g(1)=a+1$，故
\end{itemize}

\nopagebreak[4]
\[2b\geqslant-(a+1)\ \Longrightarrow\ b\geqslant-\frac{a+1}{2}.\]

\begin{itemize}
\item $a\geqslant1$ 时 $\frac1{\sqrt a}\leqslant1$，最小值为 $g\left(\frac1{\sqrt a}\right)=2\sqrt a$，故
\end{itemize}

\nopagebreak[4]
\[2b\geqslant-2\sqrt a\ \Longrightarrow\ b\geqslant-\sqrt a .\]

\gkanswer{答案：（1）$b^2>a$；（2）$0<a<1$ 时 $b\geqslant-\frac{a+1}{2}$；$a\geqslant1$ 时 $b\geqslant-\sqrt a$。}

\end{answer}
\begin{answer}
\gksolhead{22}{}

\textbf{题目}：

设 $m\in\mathbf R$，在平面直角坐标系中，已知向量 $\boldsymbol a=(mx,y+1)$，向量 $\boldsymbol b=(x,y-1)$，$\boldsymbol a\perp\boldsymbol b$，动点 $M(x,y)$ 的轨迹为 $E$．

（1）求轨迹 $E$ 的方程，并说明该方程所表示曲线的形状；

（2）已知 $m=\dfrac14$，证明：存在圆心在原点的圆，使得该圆的任意一条切线与轨迹 $E$ 恒有两个交点 $A,B$，且 $OA\perp OB$（$O$ 为坐标原点），并求出该圆的方程；

（3）已知 $m=\dfrac14$，设直线 $l$ 与圆 $C:x^2+y^2=R^2$（$1<R<2$）相切于 $A_1$，且 $l$ 与轨迹 $E$ 只有一个公共点 $B_1$，当 $R$ 为何值时，$|A_1B_1|$ 取得最大值？并求最大值．

\Needspace{4\baselineskip}
\textbf{解析}

\Needspace{5\baselineskip}
（1）由 $\boldsymbol a\perp\boldsymbol b$：

\nopagebreak[4]
\[\boldsymbol a\cdot\boldsymbol b=mx\cdot x+(y+1)(y-1)=mx^2+y^2-1=0,\]

\Needspace{5\baselineskip}
即轨迹

\nopagebreak[4]
\[E:\ mx^2+y^2=1 .\]

\begin{itemize}
\item $m=0$：$y^2=1$，即 $y=\pm1$，两条平行直线；
\item $m>0$：$\frac{x^2}{1/m}+y^2=1$，焦点在 $x$ 轴上的椭圆；
\item $m<0$：$y^2-|m|x^2=1$，焦点在 $y$ 轴上的双曲线。
\end{itemize}

（2）$m=\frac14$ 时 $E:\frac{x^2}{4}+y^2=1$，取圆 $x^2+y^2=\frac45$（半径 $\frac2{\sqrt5}$，介于短半轴 $1$ 与长半轴 $2$ 之间）。

\Needspace{5\baselineskip}
设该圆的切线 $l:y=kx+t$，则

\nopagebreak[4]
\[\frac{|t|}{\sqrt{1+k^2}}=\frac2{\sqrt5}\ \Longrightarrow\ t^2=\frac45(1+k^2).\]

\Needspace{5\baselineskip}
代入椭圆：

\nopagebreak[4]
\[\left(\frac14+k^2\right)x^2+2ktx+t^2-1=0 .\]

\Needspace{5\baselineskip}
设 $A(x_1,y_1)$、$B(x_2,y_2)$，由韦达定理

\nopagebreak[4]
\[x_1+x_2=-\frac{2kt}{k^2+\frac14},\qquad x_1x_2=\frac{t^2-1}{k^2+\frac14}.\]

\nopagebreak[4]
\[\overrightarrow{OA}\cdot\overrightarrow{OB}=x_1x_2+y_1y_2
=x_1x_2+(kx_1+t)(kx_2+t)
=(1+k^2)x_1x_2+kt(x_1+x_2)+t^2 .\]

\Needspace{5\baselineskip}
两边乘以 $k^2+\frac14$：

\nopagebreak[4]
\[(1+k^2)(t^2-1)-2k^2t^2+t^2\left(k^2+\frac14\right)
=t^2\left(1+k^2-2k^2+k^2+\frac14\right)-(1+k^2)
=\frac54t^2-(1+k^2).\]

由 $t^2=\frac45(1+k^2)$ 知上式为零，故 $\overrightarrow{OA}\cdot\overrightarrow{OB}=0$，即 $OA\perp OB$。

（切线与 $y$ 轴平行时取 $x=\pm\frac2{\sqrt5}$，由对称性同样有 $OA\perp OB$。）

\Needspace{5\baselineskip}
（3）$1<R<2$ 时，与圆 $x^2+y^2=R^2$ 相切且与椭圆只有一个公共点的直线就是椭圆的切线。设其单位法向量为 $(\cos\theta,\sin\theta)$，则

\nopagebreak[4]
\[l:\ x\cos\theta+y\sin\theta=R .\]

\Needspace{5\baselineskip}
椭圆 $\frac{x^2}{4}+y^2=1$ 沿该法向的支撑长度为 $\sqrt{4\cos^2\theta+\sin^2\theta}$，故

\nopagebreak[4]
\[R^2=4\cos^2\theta+\sin^2\theta=1+3\cos^2\theta\ \Longrightarrow\ \cos^2\theta=\frac{R^2-1}{3}.\]

\Needspace{5\baselineskip}
切点分别是

\nopagebreak[4]
\[A_1=(R\cos\theta,R\sin\theta),\qquad B_1=\left(\frac{4\cos\theta}{R},\frac{\sin\theta}{R}\right).\]

\Needspace{5\baselineskip}
于是

\nopagebreak[4]
\[|A_1B_1|^2=\cos^2\theta\left(R-\frac4R\right)^2+\sin^2\theta\left(R-\frac1R\right)^2 .\]

\Needspace{5\baselineskip}
记 $u=\cos^2\theta=\frac{R^2-1}{3}$：

\nopagebreak[4]
\[|A_1B_1|^2=u\left(R^2-8+\frac{16}{R^2}\right)+(1-u)\left(R^2-2+\frac1{R^2}\right)
=R^2-2+\frac1{R^2}+u\left(-6+\frac{15}{R^2}\right).\]

\Needspace{5\baselineskip}
代入 $u=\frac{R^2-1}{3}$：

\nopagebreak[4]
\[u\left(-6+\frac{15}{R^2}\right)=\frac{(R^2-1)(15-6R^2)}{3R^2}
=\frac{21R^2-6R^4-15}{3R^2}=7-2R^2-\frac5{R^2}.\]

\Needspace{5\baselineskip}
故

\nopagebreak[4]
\[|A_1B_1|^2=R^2-2+\frac1{R^2}+7-2R^2-\frac5{R^2}=5-\left(R^2+\frac4{R^2}\right).\]

\Needspace{5\baselineskip}
由均值不等式

\nopagebreak[4]
\[R^2+\frac4{R^2}\geqslant4,\]

\Needspace{5\baselineskip}
等号当且仅当 $R^2=2$，即 $R=\sqrt2\in(1,2)$。此时

\nopagebreak[4]
\[|A_1B_1|^2=5-4=1\ \Longrightarrow\ |A_1B_1|=1 .\]

故 $R=\sqrt2$ 时 $|A_1B_1|$ 取得最大值 $1$。

\gkanswer{答案：（1）$mx^2+y^2=1$；$m=0$ 时为两条平行直线，$m>0$ 时为椭圆，$m<0$ 时为双曲线；（2）$x^2+y^2=\frac45$；（3）$R=\sqrt2$ 时 $|A_1B_1|$ 最大，最大值为 $1$。}

\end{answer}